% ===== §11 The Dialectical-Materialist History of Intimate Power =====
\section{The Dialectical-Materialist History of Intimate Power}
\label{p17:sec:histmat}

\noindent
The two summits completed the account of power as a structure: how it freezes and must be kept from freezing, how its asymmetry can dominate or generate. But a structure is not a fate, and to leave the account there would be to present the power-relations of an intimate bond as though they were given, timeless, simply there to be managed. They are not given. They are the historical sediment of particular material conditions, and what has been sedimented can, in time, be dissolved. This section takes up the third of the paper's questions---how love becomes free within history---and it requires a method the paper has not yet deployed: a historical materialism. But it requires that method in a specific form, and the section's first task is to fix which form, because the wrong form would undo, in a sentence, everything the paper has built.

\subsection{A statement of standpoint: dialectical, not mechanical}
\label{p17:sub:histmat-standpoint}

There are two historical materialisms, and the difference between them is decisive here. The first, which may be called \emph{mechanical}, holds that the economic base unilaterally determines the superstructure: the relations of production fix the legal, ethical, symbolic, and affective forms of life, which are their passive reflection. Taken strictly, this version reduces love to an effect of its economic base---and would thereby annihilate, at a stroke, the real unsymbolisable core the paper has laboured to protect, dragging the bloom back into the order of determined conditions and calling it determined too.

The section stands, instead, on \emph{dialectical} historical materialism, and the distinction is not a softening but a fidelity. The dialectical version holds three things the mechanical one does not. First, that history advances not by smooth determination but by the motion of \emph{contradiction}---the tension between forces and relations of production, the unity and struggle of opposites---so that the engine of change is internal tension, not one-way causation. Second, that the superstructure possesses \emph{relative autonomy} and \emph{reacts back} upon the base: the symbolic, ethical, and affective life that material conditions shape is not their passive reflection but a force that in turn conditions and, under conditions, reshapes them.\footnote{The reduction of historical materialism to a one-way economic determinism was repudiated within the tradition itself; the relative autonomy of the superstructure and its reaction upon the base belong to the dialectical statement of the doctrine, not to a revision of it. The present account takes this as the orthodox rather than the heterodox reading. For the foundational analysis of how the social relations of a mode of production sediment into apparently natural forms, see \textcite{marx1976}; for the material history of the reproduction of life and its bearing on the forms of intimate and familial power, \textcite{federici2004}.} Third, that history moves by the \emph{negation of the negation}, a spiral and not a line---which is, the section will note, the very form the series has used throughout.

\begin{claim}[Constraint without determination]
On the dialectical reading, ``material conditions constrain but do not determine'' is not a dilution of historical materialism but its orthodox statement: the base exerts determining force, \emph{and} the superstructure---symbolic, ethical, affective---retains a relative autonomy through which it reacts upon the base. The mechanical reading, which makes the superstructure a passive reflection and reduces love to an effect of its economic base, is the deviation from the dialectic, not its rigour. The paper's ``safe version'' is therefore the faithful version: it preserves both the constraining force of conditions and the agency and contingency that the freezing of power would otherwise be granted as fate.
\end{claim}

\subsection{The historical genesis of power-structure}
\label{p17:sub:histmat-genesis}

With the standpoint fixed, the first substantive claim follows. The power-structure of an intimate field is not a timeless form of ``the family'' or ``the couple'' but the sediment of particular relations of production. The authority of the powerful elder is not an eternal feature of kinship; it rests on a specific economic base---the control of resources, of inheritance, of the bestowal and withdrawal of recognition---and on the superstructure that base raised up: the ethics of filial duty, the institutions of lineage, the social functions assigned to marriage. These are historical formations, and they move. As the material base shifts---as the children of a generation attain economic independence, as geographic mobility loosens the grip of locality, as the social functions once concentrated in the family migrate to other institutions---the elder's power, which rested on that base, historically wanes.

This is of the first importance for the negative summit (\xref{p17:sec:balance}), to which it supplies a material ground for hope. The freezing of power is not a metaphysical fate but the product of transient material conditions; and conditions change. Counter-force, therefore, is not the futile opposition of an eternal structure but the work---sometimes the patient awaiting, sometimes the active hastening---of a change in the material conditions on which a transient structure rests. The dialectical-materialist standpoint converts the suppression of freezing from a Sisyphean labour against the permanent into a historical labour with the grain of the possible.

\subsection{Marriage as the dialectical sublation of two material histories}
\label{p17:sub:histmat-sublation}

The section's central scene is the one that first motivated it: that a union joins not two individuals but \emph{two material histories}. A marriage, dialectically understood, is the bringing-together of two sets of relations of production, two material pasts, two ethico-symbolic systems each sedimented from a different economic base, into one contested present. Each partner carries into the new field the grammar of trust (\xref{p17:sec:encoding}), the expectations of role, the very forms of affective expression, that their own material history laid down; and these do not merely coexist. They contradict.

\begin{claim}[The union of two material histories]
A union is not the joining of two individuals but the integration of two material histories---two sets of relations of production and the ethico-symbolic systems sedimented from them. This supplies a material ground for the internal causes: the mismatch of trust-grammars (\xref{p17:sec:encoding}) is, at depth, the contradiction of two political economies of affective expression; the external power of elders (\xref{p17:sub:ir-gradient}) is a determinate historical form of the family, waning as its material base shifts; and regime bleeding (\xref{p17:sub:ir-gradient}) finds its historical explanation in a new field not yet possessed of an economic base independent of the old relations of production, and so still embedded in them.
\end{claim}

\noindent
But the bringing-together is, when it goes well, neither the victory of one material history over the other nor their inert juxtaposition. It is a \emph{sublation}---an \emph{Aufhebung}---in which the contradiction of the two systems (the moment of the negative) is worked, through the practice of the relational field, into a new and shared grammar that did not exist in either: a third thing, carrying a positive holonomy, irreducible to either of its sources. And the dialectical point is exact: sublation does not abolish the difference of the two histories but \emph{preserves} it within a higher unity---which is precisely the redefinition of connection the negative summit reached by another road (\xref{p17:sub:balance-connection}). The good union does not erase the two material histories into sameness; it lifts their contradiction into a shared form that keeps the difference alive and generative. Here the dialectical-materialist method and the holonomic account of value turn out to be describing one movement in two vocabularies: the spiral of sublation and the positive holonomy of the good cycle are the same return-with-surplus, the same negation of the negation.

\subsection{The constraint on praxis: a dialectical law of history}
\label{p17:sub:histmat-law}

The standpoint yields, finally, a constraint on the practical chapters to come, and it is the section's contribution to the praxis. The effectiveness of counter-power practice is conditioned by material circumstance. One cannot, by will or by love alone, dissolve a power-structure whose material base remains intact: where a couple's economy still depends on resources the elder can withdraw, no sufficiency of feeling will lift that power, and a counter-force mounted before its material conditions are ripe will fail---not for want of resolve but for want of a base. This is the constraint, and it disciplines the praxis against voluntarism: it explains why some resistances succeed and others fail not by the quality of their strategy alone but by the ripeness of their material moment.

But the constraint is dialectical, and so it cuts the other way as well. Material conditions \emph{open} a space of possibility; they do not fill it. Within the space the ripe conditions open, the outcome is not determined---it retains the agency of the subjects and the contingency of events. There is no necessary trajectory along which intimate liberation must run, no guarantee written into history that love will be freed. To claim otherwise would be to relapse into the teleological determinism the standpoint was chosen to refuse, and would contradict, besides, the paper's governing commitments: the far-from-equilibrium openness of the dissipative account (\xref{p17:sec:dissipation}), the irreducible remainder of the desire-structure (\xref{p17:sub:balance-desire}), and the series' standing refusal of any master-trajectory. Conditions constrain; they do not write the ending.

\begin{claim}[The dialectical law, and its two edges]
The effectiveness of counter-power practice is constrained by material conditions: a structure whose base is intact cannot be dissolved by will alone, and a counter-force mounted before its conditions are ripe fails for want of a base---which disciplines the praxis against voluntarism. But the law is dialectical and cuts both ways: ripe conditions open a space of possibility without determining its outcome, which retains the agency of subjects and the contingency of events. History constrains the liberation of love; it does not guarantee it. The first edge refutes the optimist who thinks love alone suffices; the second refutes the determinist who thinks history suffices without love.
\end{claim}

\subsection{The division of labour with the real}
\label{p17:sub:histmat-real}

One danger remains, and disposing of it welds this section to the methodological core (\xref{p17:sec:conditions}). Historical materialism, even dialectical, is a theory of how material conditions shape consciousness and its forms; does it not then threaten to reduce love itself---the real, unsymbolisable core---to a determination of the economic base, undoing the very thing \xref{p17:sec:conditions} was written to protect?

It does not, and the relative autonomy of the superstructure is precisely what prevents it. The dialectical standpoint already holds that the symbolic, ethical, and affective life is not the passive reflection of the base but reacts upon it; and the real dynamics of love are the layer of that life least reducible to material determination, the layer whose autonomy is greatest because it circles a vacancy no material condition can fill. Historical materialism, on this account, explains the \emph{conditions} of love's dynamics---the two material histories, the power-structures, the symbolic systems, the economic bases that shape the field in which love occurs---but it does not reach the dynamics themselves. It is, in the exact terms of \xref{p17:sec:conditions}, the \emph{diachronic dimension} of the political economy of conditions: where \xref{p17:sec:conditions} described the synchronic structure of the fence (symbolic means, real end), this section describes the fence's historical genesis (how the material past sediments into the present structure). The bloom remains what material conditions make room for and ward, but do not cause.

\begin{claim}[Historical materialism as the diachronic political economy of conditions]
Dialectical historical materialism explains the \emph{conditions} of love's dynamics---the sedimented material histories, power-structures, and symbolic systems that shape the field---but, by the relative autonomy of the superstructure, does not reach the real dynamics themselves. It is the diachronic dimension of the political economy of conditions: \xref{p17:sec:conditions} gave the synchronic structure of the fence, this section its historical genesis. Material conditions make room for the bloom and ward it; they do not cause it. Historical materialism is thus the temporal ally of the fence-and-bloom thesis, not its reductive rival---and the ``condition'' so gains a temporal depth (history) to set beside its spatial depth (the gradient of \xref{p17:sub:ir-gradient}), while the real remains, in both, what conditions protect but never produce.
\end{claim}

\noindent
The same dialectical method governs the series' treatment of justice elsewhere, where a juridical form adapted to an old material base enters into contradiction with new conditions, and a new form is seen struggling to be born within the contradiction rather than being deduced from the base.\footnote{For the dialectical-materialist treatment of juridical and institutional form---old form fitted to old base, contradiction sharpening, new form germinating within the contradiction rather than read off from the base---see the corresponding development in \textcite{paperxiii}.} The method is one across the series, and its lesson for the third question is now complete: love becomes free within history not by escaping material conditions, which is impossible, nor by being delivered by them, which is determinism, but by the labour of subjects working with the grain of changing conditions to dissolve the structures those conditions had sedimented---a freedom that is neither given nor guaranteed, but won, dialectically and in time. With this the three questions have their three answers, and the paper turns to set estrangement within the widest frame it commands: the cosmology of generative systems.
